\documentclass[11pt,a4paper]{article}
\usepackage[top=2.5cm,bottom=2.8cm,left=2.5cm,right=2.5cm]{geometry}
\usepackage{amsmath,amssymb,amsthm}
\usepackage{url}
\usepackage{hyperref,xcolor,enumitem,booktabs,array}
\hypersetup{colorlinks=true,linkcolor={blue!60!black},citecolor={blue!60!black},urlcolor={blue!60!black}}

\newtheorem{theorem}{Theorem}[section]
\newtheorem{proposition}[theorem]{Proposition}
\newtheorem{remark}[theorem]{Remark}
\theoremstyle{definition}\newtheorem{definition}[theorem]{Definition}

\newcommand{\vp}{\varphi}
\newcommand{\Lcond}{\Lambda_{\mathrm{cond}}}
\newcommand{\SU}{\mathrm{SU}}
\newcommand{\twoT}{2T}
\newcommand{\twoD}{2D}
\newcommand{\twoI}{2I}
\newcommand{\QH}{Q_H}
\newcommand{\Ztwo}{\mathbb{Z}_2}
\newcommand{\Zthree}{\mathbb{Z}_3}
\newcommand{\Odm}{\Omega_{\mathrm{DM}}}
\newcommand{\Ob}{\Omega_{b}}

\title{\textbf{Dark Genesis: Co-genesis of Baryons and a Sterile Dirac Dark Sector
at the First-Order Confinement Transition of the Density-Feedback Hopfion}}
\author{Frederick Manfredi}
\date{2026}

\begin{document}
\maketitle

\begin{abstract}
Paper~XXI identifies the beyond--Standard-Model D-type sector as a sterile, Dirac
dark-matter candidate. This companion develops the \emph{genesis
dynamics}: how baryons and the dark sector are produced together in the early universe.
The thread: (i)~the confinement transition is \emph{first-order}, derived from the condensate thermal free
energy of the QCD epoch: the trefoil's three colours supply the $\Zthree$ cubic of a
Polyakov-loop potential whose coefficients follow from the semi-Dirac excitation gas, and the
full free energy $F(T,L)$ shows the order parameter jump discontinuously from the confined
branch to the deconfined value at $T_c\sim\Lcond(\mathrm{QCD})$; (ii)~a chiral, CP-violating
reconnection at the bubble walls realises all three Sakharov conditions (baryon-number violation
= sector-changing reconnection, CP = the derived Wess--Zumino--Witten phase
$\delta_{\rm CP}=65.0^\circ$, departure from equilibrium = the first-order transition), with
confinement itself removing the free strands that reconnection requires, so baryon-number
violation shuts off in the broken phase automatically; (iii)~because the CP phase is a property
of the shared $\twoI$ vacuum, the same event seeds \emph{correlated} baryon and dark
asymmetries (co-genesis), and the Dirac D-type sector carries its asymmetry, giving an
asymmetric-dark-matter abundance $\Odm/\Ob\approx5$ for $m_{\mathrm{dark}}\approx5$--$6$~GeV. All
three Sakharov conditions are thus derived. The one quantity not derived is the baryon
asymmetry's \emph{magnitude}: the maximal efficiency overproduces, and the suppression that
matches $\eta\approx6\times10^{-10}$ is supplied by the heavy-band anisotropy
$S_{\rm eff}\propto\sin^4\theta/\vp^6$ of the baryon channel, controlled by the same
density-feedback combination $\beta_*\rho$ that drives the first-order transition; a
self-consistent value exists inside the first-order window but is not forced, leaving the
window's location in physical units and the heavy-band power as the genuine openings. The
co-genesis number-asymmetry ratio $r=n_{\rm dark}/n_b\sim O(1)$ is independent of all this
(the shared transport efficiency cancels), so with the canonical asymmetric-dark-matter scale
$m_{\mathrm{dark}}\sim5$~GeV it gives $\Odm/\Ob\approx5$ --- a conditional consistency that
fixes $m_{\mathrm{dark}}$, not a parameter-free prediction. We state precisely what is derived
and what remains.
\end{abstract}

\tableofcontents

\section{Scope}
Paper~XXI~\cite{Paper21} is the \emph{structural} sector paper (the sterile Dirac D-type
dark state, the no-$E_7$ theorem). The present paper is the \emph{dynamical} companion: the
production mechanism. Sections~\ref{P22:sec:fo}--\ref{P22:sec:abund} derive the mechanism --- the
first-order transition from the epoch's thermal free energy, the three Sakharov conditions, and
the co-genesis ratio. Section~\ref{P22:sec:twoepoch} then places the resulting relic in its
cosmological role --- cold dark matter at recombination and in clusters, alongside the
baryon-sourced director strain that governs galaxies --- and Section~\ref{P22:sec:open} states
the one quantity that remains open: the \emph{magnitude} of $\eta$, which reduces to the
density-feedback strength $\beta_*\rho$ at the QCD epoch, co-constrained by the requirement of a
first-order transition.

\section{The confinement transition is first-order}
\label{P22:sec:fo}
The deconfinement free energy is a Polyakov-loop Landau--Ginzburg potential for the order
parameter $L$, constrained by a $\Zthree$ center symmetry~\cite{Pisarski2000}:
\begin{equation}
  V(L)=\tfrac{b_2(T)}{2}|L|^2-\tfrac{b_3}{3}(L^3+L^{*3})+\tfrac{b_4}{4}|L|^4-h(L+L^*).
  \label{P22:eq:polyakov}
\end{equation}
The cubic $b_3$ (allowed only by $\Zthree$) drives a first-order transition; the explicit
breaking $h$ weakens it, and above a critical $h_c$ the transition becomes a crossover
(the physical light-quark point of real QCD).

\begin{proposition}[Framework first-order]
\label{P22:prop:fo}
In the density-feedback Hopfion the $\Zthree$ is the colour symmetry of the trefoil
$T_{2,3}$ (the $\Zthree$ outer automorphism of $\twoT$, Paper~XV~\cite{Paper15}), so the
cubic $b_3\neq0$ is present by construction (the colour-singlet baryon is the $\Zthree$
invariant $3\times3\times3\to1$). The explicit breaking is small, $h\sim e^{-E_{\rm sb}/T}$
with $E_{\rm sb}\sim2m_{\rm const}\sim600$~MeV the string-breaking energy and
$T\sim150$~MeV, giving $h\sim0.02$--$0.14$, below the critical endpoint $h_c\approx0.21$.
The transition is therefore first-order --- the framework differs from real QCD (crossover)
precisely in having no free light quarks: $\Zthree$ is broken only by string-breaking, not by
light dynamical quarks.
\end{proposition}
\noindent A density-feedback contribution (the $(1+\beta_*\rho)$ softening) adds to the
coefficients and provides a second, independent first-order driver (\path{firstorder_transition.py};
the first-order region is an intermediate window in $\beta_*\rho$, Section~\ref{P22:sec:open}).

\subsection{The coefficients from the geometry of the epoch}
\label{P22:subsec:coeffs}
The coefficients of Eq.~\eqref{P22:eq:polyakov} are not free: they are the holonomy expansion of
the thermal free energy of the condensate excitations, and those excitations carry the
\emph{semi-Dirac} dispersion of the medium (Paper~VII~\cite{Paper7}),
\begin{equation}
  E(k,\theta)^2=A^2k^4\sin^4\theta+v^2k^2\cos^2\theta ,
  \label{P22:eq:semidirac}
\end{equation}
linear along the easy (radial) axis and quadratic along the hard (perpendicular) axis. The
anisotropic phase space changes the temperature power: $\int\!d^3k\,e^{-nE/T}\propto(T/n)^{2}$
rather than $(T/n)^3$, so the single-phase holonomy potential is
\begin{equation}
  f(T,\chi)=-\sum_{n\ge1}c_n\cos n\chi,\qquad c_n\propto\frac{T^3}{n^{3}} ,
  \label{P22:eq:weiss}
\end{equation}
a cubic-Bernoulli ($B_3$) Weiss potential, in contrast to the relativistic quartic ($B_4$,
$c_n\propto T^4/n^4$). Summing the three $\SU(3)$ colour (holonomy) phases $(\chi,0,-\chi)$ along
the real axis, with $L=(1+2\cos\chi)/3$ and $\cos n\chi=T_n\!\big(\tfrac{3L-1}{2}\big)$ (Chebyshev), and
expanding in $L$ gives
\begin{equation}
  \frac{b_3}{b_4}\approx0.17 \quad(\text{semi-Dirac }B_3),\qquad\text{vs}\quad 0.16\ (B_4),
  \label{P22:eq:b3b4}
\end{equation}
so the epoch's geometry modestly \emph{enhances} the cubic (strengthens first-order) relative
to a relativistic gas. Confinement gaps the excitations through a $\Zthree$-even coupling to
the order parameter, $\Delta(L)^2=\Delta_0^2(1-L^2)$ with $\Delta_0\sim\Lcond(\mathrm{QCD})$;
this contributes a thermal mass $b_2(T)$ whose sign changes at $T_c\sim2\Lcond\sim T_{\rm QCD}$,
matching the framework's near-critical condition $T/\Lcond\approx2$ (Paper~VII~\cite{Paper7}).
(The three steps are computed in \path{semidirac_thermal.py} (the $T^3$/$B_3$ Weiss potential),
\path{su3_colorsum_coeffs.py} (the colour-sum $b_3/b_4$), and \path{gap_b2T.py} (the thermal
mass).)

\begin{theorem}[Derived first-order transition]
\label{P22:thm:fo}
The full thermal free energy combining the holonomy and the gap in one normalisation,
\begin{equation}
  F(T,L)=-T\sum_{n\ge1}\frac1n\,I_n\!\big(T,\Delta(L)\big)\,[\,2\cos n\chi+1\,]+a_{\rm vac}L^2,
  \qquad I_n=\int\!\frac{d^3k}{(2\pi)^3}\,e^{-nE/T},
  \label{P22:eq:FTL}
\end{equation}
has, below $T_c$, its global minimum on the confined branch ($L\!\approx\!0.35$--$0.45$) and,
above $T_c$, at the deconfined value $L\!\to\!1$, with the two degenerate at
$T_c\sim(0.9$--$1.5)\,\Delta_0\sim O(\Lcond)$. The order parameter therefore jumps
\emph{discontinuously} by $\Delta L\approx0.55$--$0.65$ across a finite barrier: the transition
is first-order, driven by the $\Zthree$ cubic (the $n=3k$ holonomy harmonics, i.e.\ the trefoil's
three-fold colour structure). (Computed in \path{unified_FTL.py}.)
\end{theorem}
\noindent The jump in the order parameter is large (strongly first-order in order-parameter
space), while the latent heat is a modest fraction of the radiation bath,
$\alpha=L_{\rm heat}/\rho_{\rm rad}\approx0.03$--$0.06$ (the barrier scale sits well below
$g_*T^4$), giving bubble-wall speeds $v_w\sim0.15$--$0.2$ (deflagration). What matters for
baryogenesis is not $\alpha$ but the departure from equilibrium at the wall and the shutoff of
baryon-number violation behind it; the latter is automatic here (Section~\ref{P22:sec:baryo}).

\section{Baryogenesis as chiral reconnection}
\label{P22:sec:baryo}
A baryon is a chiral wound state: the trefoil (writhe $+3$) differs from its mirror
(Paper~XVIII~\cite{Paper18}). The three Sakharov conditions map onto knot dynamics:
baryon-number violation is the reconnection $\QH{=}2+\QH{=}2\to\QH{=}3+\QH{=}1$ (baryons form
from leptons by threading); C and CP violation are the derived Wess--Zumino--Witten phase
$\delta_{\rm CP}=65.0^\circ$ (Paper~V~\cite{Paper5}) together with the chiral writhe; and the
departure from equilibrium is the first-order transition of Section~\ref{P22:sec:fo}. Because the
reconnection rate is a strong-interaction rate, $\Gamma/H\sim10^{19}$ at $T_{\rm QCD}$ (deep
equilibrium); only a genuine first-order freeze-out escapes the resulting washout. Here that
escape is structural: a reconnection needs free trefoil strands, which exist only in the
deconfined phase, so once the bubble of confined phase forms the baryon-number-violating rate
shuts off \emph{automatically} --- the analogue of the electroweak ``sphaleron shutoff in the
broken phase'' is guaranteed by confinement itself, with no condition on the transition
strength. Since $\sin\delta_{\rm CP}=0.91$ is of order unity, the CP bias is not the limiting
factor.

\section{Co-genesis and the asymmetric-dark-matter abundance}
\label{P22:sec:abund}
The phase $\delta_{\rm CP}$ is a property of the shared $\twoI$ vacuum, which contains both
$\twoT$ (baryon) and the D-type $\twoD_n$ (dark) as subgroups. The same chiral, CP-violating
reconnection at the bubble walls therefore biases \emph{both} freeze-outs, seeding correlated
asymmetries. The D-type sector is Dirac (its odd-rank spinor is complex, Paper~XXI), so it
carries a dark number and an asymmetry; its symmetric part self-annihilates (the sector is
self-interacting through the shared director field), leaving the asymmetric residue. Its
abundance follows the asymmetric-dark-matter arithmetic,
\begin{equation}
  \frac{\Odm}{\Ob}=r\,\frac{m_{\rm dark}}{m_p},\qquad r=\frac{n_{\rm dark}}{n_b},
  \label{P22:eq:adm}
\end{equation}
with $r$ the ratio of the dark and baryon number asymmetries.

One feature of this is robust: $r$ is \emph{independent of the unpinned magnitude} of $\eta$
(Section~\ref{P22:sec:open}), because both asymmetries are seeded by the same $\delta_{\rm CP}$ at the
same transition, so the common transport efficiency cancels in the ratio. What remains \emph{not}
derived is the value of either $r$ or $m_{\rm dark}$. The branching is set at the shared freeze-out
by each sector's thermal degrees of freedom, which for a Wess--Zumino--Witten sector is its central
charge $c$; with the common $\delta_{\rm CP}$ bias cancelling,
$r\simeq c_{\rm dark}/c_{\rm baryon}=c(\mathrm{SO}(10)_1)/c((E_6)_1)=5/6\approx0.83$ for $\twoD_3$
(and $7/6$ for $\twoD_5$), consistent with the cruder subgroup-order ($12/24$) and primary counts
and so robustly $O(1)$. This is a thermal-DOF estimate, not yet a full reconnection-amplitude rate:
the residual is the per-channel CP-asymmetry ratio, here taken as unity. Because the D-type carries
no $Q_{\rm group}$ (Paper~XXI), its mass tower is undefined and the framework does \emph{not} pin
$m_{\rm dark}$; with $r\simeq0.83$, matching the observed $\Odm/\Ob\approx5.4$ then \emph{fixes}
$m_{\rm dark}\approx6$~GeV --- the canonical asymmetric-dark-matter scale, a sterile $\sim$few-GeV
state that evades direct detection by construction. This is therefore a \emph{conditional
consistency}, not a parameter-free prediction: it rests on the thermal-DOF weight $r\approx0.8$
(its per-channel amplitude correction uncomputed), on efficient symmetric self-annihilation through the
sector's SIDM channel, and on the dark relic \emph{existing} --- the open question of
Paper~XXI~\cite{Paper21}; its \emph{stability}, once it exists, follows there from the
spinor ($\mathrm{SO}(10)$ matter) parity.

\section{The D-type sector as two-epoch dark matter}
\label{P22:sec:twoepoch}
The D-type relic and the baryon-sourced director \emph{strain} of
Paper~VII~\cite{Paper7} are two configurations of the one condensate, and they
dominate the dark-matter phenomenology at two different epochs --- a single fabric
in two modes, not two dark species.

\paragraph{Recombination and clusters: the relic as cold dark matter.}
Paper~VII identifies the \emph{galactic} dark matter as the director strain, but that
component is late and isocurvature: its ultralight scale $m_\xi\sim H_0$ keeps it
Hubble-frozen until $z\sim0$, so it cannot supply the cold matter the acoustic peaks
require at recombination. The D-type relic is precisely the missing ingredient ---
non-baryonic, electromagnetically and weakly sterile (Paper~XXI~\cite{Paper21}),
non-ultralight ($m_{\rm dark}\sim5$~GeV$\,\gg H(z_{\rm rec})$), cold, and, being seeded
by the co-genesis of Section~\ref{P22:sec:abund} correlated with the baryon asymmetry,
\emph{adiabatic} rather than isocurvature. It therefore clusters as standard cold dark
matter at recombination, and at the cosmic ratio $\Odm/\Ob\approx5$ supplies the cluster
lensing mass (collisionless, as in $\Lambda$CDM) where the baryon-tracking strain alone
falls short. This collisionless relic is the merger-\emph{independent} bulk of that mass;
the baryon-independent disclination network of Paper~VII~\cite{Paper7} is then not a
competing account of the same deficit but the merger-\emph{dependent} modulation riding on
top of it --- denser, and so a larger excess, in recently-disturbed clusters --- which is
where the sector's distinctive non-$\Lambda$CDM, non-MOND signature lies.

\paragraph{Galaxies: the strain, and why the relic does not spoil the RAR.}
At $z\sim0$ the strain unfreezes and yields the tight radial-acceleration relation (RAR)
because it is sourced by and tracks the baryons (Paper~VII~\cite{Paper7}). A cosmic-ratio
relic, were it collisionless here too, would build halos that do not track baryons and
would spoil that tightness. It does not, because the relic self-interacts through the
\emph{shared} director field --- the single $\hat n$ that carries the Bell correlations of
Papers~VIII--IX~\cite{Paper8,Paper9} --- and that interaction inherits the field's
anisotropy. The mediator is the semi-Dirac director mode $E^2=v^2k_\parallel^2+A^2k_\perp^4$,
and the suppression $S_{\rm eff}=\sin^4\theta/[\vp^6(1+\beta_*\rho)]$ is the \emph{same}
kernel that fixes those correlations; its density factor $(1+\beta_*\rho)$ is negligible at
every diffuse dark-matter density, screening only inside stars and the knot cores. The
decisive feature is the geometry-dependence this kernel enforces, established for the Bell
observable in Papers~VIII--IX: $\sin^4\theta$ is \emph{resolved} in a coherent geometry but
\emph{washes out} under isotropic averaging (planar Bell tests reach
CHSH$\,\approx2.6$--$2\sqrt2$, a fully isotropic $3$D geometry falls to the classical range).
The relic interaction inherits this exactly. Where the baryon-sourced strain makes $\hat n$
\emph{coherent}, inside a galaxy, the interaction is resolved and active: it relaxes the
relic, which thermalises to the baryon-dominated potential, tracks the baryons, and
preserves the RAR --- plasma-like, $t_{\rm relax}\sim v^3/(nC^2)$, with a weak coupling
$C\sim10$--$10^3$~eV. Where $\hat n$ is \emph{isotropic} --- the smooth recombination
medium, and the inter-galactic field --- the interaction washes out and is inert, so the
relic is collisionless cold dark matter there, with no long-range force and no cosmological
over-closure. The galaxy/cluster split and the recombination safety are thus the same
geometry-gating the Bell sector already exhibits, not an added assumption; the resulting
dark sound speed in the coherent regime, $c_s^2=C/m_{\rm dark}\sim10^{-8}$, keeps the Jeans
scale far below the acoustic-peak scales. This same separation of speeds --- the lensing
potential tracks the fast-settled geometry while the director strain settles only at
$c_{\rm dir}\to0$ --- is why a galaxy's weak-lensing mass and its kinematic (dynamical) mass
need not coincide: the lens reads the settled geometry, the dynamics lag, and the two differ
at $O((c_{\rm dir}/v)^2)$ (Paper~VII~\cite{Paper7}).

\paragraph{Status.} This reproduces the observed division --- $\Lambda$CDM-like at the
microwave background and in clusters, a MOND-like RAR in galaxies --- from two modes of one
condensate, distinct from either pure theory, and it ties the dark-matter self-interaction
to the \emph{same} geometry-dependent $\sin^4\theta$ structure as the Bell correlations
(Papers~VIII--IX~\cite{Paper8,Paper9}): one anisotropy, two arenas. It is conditional on the
relic's existence and abundance (Section~\ref{P22:sec:abund}; Paper~XXI~\cite{Paper21}). The
cosmological over-closure a long-range logarithmic force would otherwise carry is removed by
the isotropic washing-out above. Within a coherent halo the interaction is active but does
\emph{not} over-close: the mean distortion is a single topological disclination (energy
$\propto R$ --- the halo mass itself --- not $\propto N^2$), the relic--relic fluctuations
Debye-screen in the thermalised plasma (the screening length falls far below the halo radius
for the coupling window above, making the correlation energy extensive, $\propto N$), and the
marginally-confining logarithm virialises rather than runs away; no further nonlinear ($S^2$)
saturation is required. The residual assumption is the quasi-neutrality of those fluctuations
about the mean texture, and the coupling $C$ (hence the screening length) is cutoff-set. Shared, cross-domain
falsifiable signatures: the planar-versus-spherical Bell geometry-dependence (testable with
$3$D or satellite entanglement protocols, Papers~VIII--IX) probes the same kernel that gates
the dark sector; a sterile $\sim$few-GeV self-interacting relic; a mild suppression of
sub-megaparsec power; and weak lensing exceeding the kinematic mass in galaxies by
$O((c_{\rm dir}/v)^2)$.

\section{Open problem: the magnitude of $\eta$}
\label{P22:sec:open}
All three Sakharov conditions are derived (Sections~\ref{P22:sec:fo}--\ref{P22:sec:baryo}): the
first-order transition from the epoch's thermal free energy (Theorem~\ref{P22:thm:fo}), the
order-unity CP phase $\delta_{\rm CP}=65.0^\circ$, and the automatic shutoff of baryon-number
violation in the confined phase. What is \emph{not} derived is the \emph{magnitude} of the
baryon asymmetry $\eta=7.04\,n_B/s$. We write
\begin{equation}
  \frac{n_B}{s}\sim A_{\rm pref}\,\varepsilon_{\rm CP}\,\delta_{\rm neq}\,b_{\rm ch},\qquad
  A_{\rm pref}=\frac{15}{4\pi^2 g_*}\sim2\times10^{-2},\qquad
  \varepsilon_{\rm CP}=\sin\delta_{\rm CP}\,c_{\rm int},
  \label{P22:eq:eta}
\end{equation}
and examine each factor. Two are \emph{not} small. By CPT the asymmetry vanishes in exact
equilibrium, and the reconnection rate is a strong-interaction rate ($\Gamma/H\sim10^{19}$), so
the survivor is the departure from equilibrium at the wall, $\delta_{\rm neq}=v_w/(\Gamma L_w)
\sim10^{-2}$ (set by the derived $v_w\sim0.2$). The CP interference is likewise $O(1)$: the
phase of Section~\ref{P22:sec:baryo} is the argument of a \emph{two-amplitude} sum (Paper~V%
~\cite{Paper5}) whose contributions are comparable, giving an interference visibility
$c_{\rm int}\sim0.6$. With a naive isotropic channel fraction $b_{\rm ch}\sim0.2$, these factors
\emph{overproduce} $\eta$ by $\sim10^{4}$--$10^{5}$: in this framework reconnection baryogenesis
is efficient, so $\eta$'s problem is a \emph{suppression}, not a smallness. (The factor anatomy
is \path{transport_eta.py} and \path{branching_eta.py}; the CP visibility uses the two-amplitude
structure of Paper~V.)

The suppression is supplied by the \emph{anisotropy}, and only appears once the directional
structure of the medium is resolved rather than angle-averaged. The baryon ($(E_6)_1$, $Q_H{=}3$)
is the \emph{heavy band} of the semi-Dirac dispersion, and the heavy band carries the directional
suppression of the condensate (Paper~VII~\cite{Paper7}, Axiom~A2),
\begin{equation}
  S_{\rm eff}=\frac{\sin^4\theta}{\vp^6\,(1+\beta_*\rho)},
  \label{P22:eq:seff}
\end{equation}
whereas the competing light-band channels (the elastic channel, the $Q_H{=}1$ neutrino byproduct)
do not. The baryon-forming reconnection therefore picks up $b_{\rm ch}\to b_{\rm ch}\,S_{\rm eff}^{\,p}$.
The power is fixed by the final state: the reconnection $2+2\to3+1$ produces exactly one heavy-band
object (the $Q_H{=}3$ baryon; the $Q_H{=}1$ neutrino is light-band), so the amplitude carries one
factor of $S_{\rm eff}$, and $\eta$, a rate $\propto|\mathrm{amplitude}|^2$, carries $S_{\rm eff}^2$:
$p=2$ (Paper~XXI~\cite{Paper21}). The factor $\vp^6$ (the golden
ratio) is the Wess--Zumino--Witten / Bogomolny power (a \emph{physical}, not frame-dependent,
power), so this is a real suppression of $\sim\vp^{-6}\approx0.056$ per power.

Crucially, the strength $\beta_*\rho$ in Eq.~\eqref{P22:eq:seff} is the \emph{same} density-feedback
combination that drives the first-order transition (Section~\ref{P22:sec:fo}, the $(1+\beta_*\rho)$
softening). It is one variable doing two jobs. Because the condensate density $\rho$ tracks $T^4$,
$\beta_*\rho$ is large at the dense QCD epoch --- which is exactly the source of the baryon-channel
suppression --- but it cannot be arbitrarily large: deep saturation softens the barrier away and
returns a crossover, so a first-order transition exists only in an \emph{intermediate window} of
$\beta_*\rho$. Overlaying the two constraints on the single axis $\beta_*\rho$, at the pinned $p=2$: the
self-consistent value $\beta_*\rho\sim15$ lies \emph{inside} the first-order window and gives
$\eta\approx6\times10^{-10}$, with the window bracketing the observed value (overproducing at
smaller $\beta_*\rho$, underproducing at larger). This $\sim15$ is the value the saturated trefoil
configuration must return (below); it is not an independent input and, as we show below, is bounded
to the right range by the established trefoil sharpness. (The anisotropy-resolved branching is \path{branching_anisotropic.py}; the
overlay of the first-order window and the $\eta$ constraint on the single axis $\beta_*\rho$ is
\path{beta_rho_window.py}.)

The magnitude of $\eta$ has therefore reduced from a transport unknown spanning many orders to a
single density-feedback quantity, $\beta_*\rho$ at the QCD epoch, co-constrained by the
requirement of a first-order transition and shared with the density screening of the dark sector.
Most of the suppression is moreover already \emph{fixed}: the per-power factor
$1/[\vp^6(1+\beta_*\rho)]$ carries the physical (Bogomolny) $\vp^6$, so for $p=2$ the
$\vp^{12}\approx322$ supplies all but a residual factor $\sim2.5\times10^2$, which is the feedback
term $(1+\beta_*\rho)^2$ with $\beta_*\rho\sim15$. That the feedback is \emph{saturated}
($\beta_*\rho\gg1$) here is expected for a sharp reason: baryogenesis \emph{is} trefoil formation
(the reconnection $2+2\to3+1$ assembles the baryon), so the feedback acting on the baryon-forming
channel is the baryon (trefoil) sector's \emph{own} --- the very configuration whose feedback
saturates and denies the quark a static saddle (Paper~XX~\cite{Paper20}). The gradient reading there
($\beta g_2$, $g_2=|\nabla n|^2$) and the density reading here ($\beta_*\rho$) are two descriptions
of the \emph{same} trefoil texture, related by the condensate's $\rho\propto\kappa$
(Paper~VII~\cite{Paper7}); they are one saturation, not a coincidence of two, and $\beta_*$ is the
density-modified, self-consistent coupling of that configuration, not the vacuum constant. What this
does \emph{not} yet fix is the \emph{magnitude}: the self-consistent trefoil value of $\beta_*\rho$
is determined in principle but its extraction is gated by the trefoil resolution wall, and the
$\rho\propto\kappa$ relation carries an as-yet-unpinned constant. An order-of-magnitude plausibility
follows from the \emph{established} trefoil sharpness --- core radius $\sim\!1/C^\ast$,
$C_3^\ast=2.5062$ (Paper~XV~\cite{Paper15}), giving $g_2\sim C^{\ast2}\approx6$ and a core feedback of
order a few --- which places $\beta_*\rho$ in the saturated range rather than at a free value. We do
not claim more: a direct evaluation of the self-consistent core feedback on a charge-resolving
trefoil ansatz is required to turn this into a number, and remains open, both because the sharp tube
under-resolves the Hopf charge at accessible grids and because the numeric feedback must be
calibrated to the $S_{\rm eff}$ normalisation used here. So $\beta_*\rho\sim15$ is the value the
saturated trefoil configuration must return, not an independent input. With the power now fixed
($p=2$, above) the single remaining open quantity is the \emph{precise} value of $\beta_*\rho$, and
whether the self-consistent trefoil core in fact reaches $\sim15$ --- order-of-magnitude plausible
but not yet established.
The co-genesis ratio $\Odm/\Ob\sim O(5)$ (Section~\ref{P22:sec:abund}) is
independent of all of this and stands as the sector's falsifiable prediction.

\begin{thebibliography}{99}
\bibitem{Paper5} F.~Manfredi, \textit{Colour Confinement, the QCD Scale, and Hadronic Structure
from the Density-Feedback Hopfion}, Zenodo (2026).
\href{https://doi.org/10.5281/zenodo.19638857}{doi:10.5281/zenodo.19638857}

\bibitem{Paper7} F.~Manfredi, \textit{Emergent Gravity, Inflation, Dark Energy, and the Weak
Equivalence Principle from the Density-Feedback Hopfion Condensate}, Zenodo (2026).
\href{https://doi.org/10.5281/zenodo.19646953}{doi:10.5281/zenodo.19646953}

\bibitem{Paper8} F.~Manfredi, \textit{Geometry-Dependent Bell Violations from the
Density-Feedback Hopfion}, Zenodo (2026).
\href{https://doi.org/10.5281/zenodo.18737070}{doi:10.5281/zenodo.18737070}

\bibitem{Paper9} F.~Manfredi, \textit{Division Algebras, the Born Rule, and the Tsirelson
Bound from the Density-Feedback Hopfion}, Zenodo (2026).
\href{https://doi.org/10.5281/zenodo.20075227}{doi:10.5281/zenodo.20075227}

\bibitem{Paper15} F.~Manfredi, \textit{The $Q_H=3$ Sector of the Density-Feedback Hopfion},
Zenodo (2026). \href{https://doi.org/10.5281/zenodo.20691001}{doi:10.5281/zenodo.20691001}

\bibitem{Paper18} F.~Manfredi, \textit{Preimage Topology and Confinement: The $Q_H=3$ Sector
of the Density-Feedback Hopfion}, Zenodo (2026).
\href{https://doi.org/10.5281/zenodo.21047750}{doi:10.5281/zenodo.21047750}

\bibitem{Paper20} F.~Manfredi, \textit{The Dynamical Layer of the Quark Sector: $E_6$-Native
Generations, the Isospin Mass-Scale, and Static Residuals as Strong Dynamics}, Zenodo (2026).
\href{https://doi.org/10.5281/zenodo.23047960}{doi:10.5281/zenodo.23047960}

\bibitem{Paper21} F.~Manfredi, \textit{The Beyond--Standard-Model D-type Sector of the
Density-Feedback Hopfion: A Sterile, Dirac Dark-Matter Candidate}, Zenodo (2026).
\href{https://doi.org/10.5281/zenodo.23319530}{doi:10.5281/zenodo.23319530}

\bibitem{Pisarski2000} R.~D.~Pisarski, \textit{Quark-gluon plasma as a condensate of
$\SU(3)$ Wilson lines}, Phys.\ Rev.\ D \textbf{62}, 111501 (2000).
\href{https://doi.org/10.1103/PhysRevD.62.111501}{doi:10.1103/PhysRevD.62.111501}

\end{thebibliography}

\end{document}
